\par\medskip\noindent
\underline{3.1(c).} Si $B$ es un objeto de
$\underline C_{A|R}$, el diagrama $(*)$ de \sisi{2.1(a)}{3.1(a)}, junto con la
exactitud de $\widetilde j$ y $j^{\vee}$, proporciona un isomorfismo
canónico
\[
 H^{\bullet}(B|R,\widetilde jF)=
 [\underline H^{\bullet}_{A|R}(F)](B).
\]
Por otra parte, un homomorfismo de anillos $S\longrightarrow R$ induce el
diagrama $(**)$ de \sisi{2.1(a)}{3.1(a)}. El funtor $\widetilde\beta$ no tiene por qué
ser exacto; sin embargo, $\beta$ conserva los productos fibrados, y se
obtiene una sucesión espectral
\[
 \underline H^p_{A|R}(R^q\widetilde\beta F)
 \Longrightarrow
 \beta^{\vee}\underline H^{p+q}_{A|S}(F),
\]
o, al evaluar en $A$,
\[
 H^p(A|R,R^q\widetilde\beta F)
 \Longrightarrow H^{p+q}(A|S,F).
\]
Como hay un isomorfismo canónico
\[
 \widetilde\beta=
 G_{A|R}\circ\beta^{\vee}\circ i_{A|S}
\]
y el funtor
$\beta^{\#}=G_{A|R}\circ\beta^{\vee}$ es exacto, se tiene

%%%%%%%%%%  scan idx 85  |  folio 74  %%%%%%%%%%

\[
 R^{\bullet}\widetilde\beta F=
 \beta^{\#}\underline H^{\bullet}_{A|S}(F).
\]
Así, la sucesión espectral precedente puede escribirse
\[
 H^p\bigl(A|R,\beta^{\#}\underline H^q_{A|S}(F)\bigr)
 \Longrightarrow H^{p+q}(A|S,F).
\]

\underline{Lema 3.2.} Sea
$B\in\operatorname{Ob}(\underline C_{A|R})$, y elíjase una flecha
suprayectiva $P\longrightarrow B$ en $\underline C_{B|R}$, donde $P$ es
un álgebra polinómica sobre $R$. Entonces, para todo
$F\in\operatorname{Ob}(\overset{\vee}{\underline C}_{A|R})$, la
aplicación natural
\[
 H^{\bullet}(\{P\longrightarrow B\},F)
 \longrightarrow[\check{\underline H}^{\,\bullet}(F)](B)
\]
es un isomorfismo.

\underline{Demostración.} Basta demostrar que la aplicación natural
\[
 H^0(\{P\longrightarrow B\},F)
 \longrightarrow[\check{\underline H}^{\,0}(F)](B)
 =\varinjlim_{\{B'\to B\}\in\operatorname{Cov}(B)}
   H^0(\{B'\longrightarrow B\},F)
\]
es un isomorfismo. Esto es claro porque, para todo
$\{B'\longrightarrow B\}\in\operatorname{Cov}(B)$, hay una aplicación
de recubrimientos
\[
 \{P\longrightarrow B\}\longrightarrow
 \{B'\longrightarrow B\}.
\]
\hfill$\square$

\underline{Corolario 3.3.}

\begin{enumerate}[label=(\alph*)]
\item Supongamos que $A$ es noetheriano. Si
$F\in\operatorname{Ob}(\widetilde{\underline C}_{A|R})$ es de tipo
finito, entonces $\underline H^n(F)$ es de tipo finito para todo $n$.

\item Si $P$ es un álgebra polinómica sobre $R$, entonces
\[
 H^n(P|R,F)=0
\]
para todo $n>0$ y todo
$F\in\operatorname{Ob}(\widetilde{\underline C}_{P|R})$.

\item Una sucesión
\[
 0\longrightarrow F'\longrightarrow F\longrightarrow F''
 \longrightarrow0
\]
en $\widetilde{\underline C}_{A|R}$ es exacta si y solo si
\[
 0\longrightarrow F'(P)\longrightarrow F(P)
 \longrightarrow F''(P)\longrightarrow0
\]
es exacta para toda $R$-álgebra polinómica
$P\in\operatorname{Ob}(\underline C_{A|R})$.
\end{enumerate}

\underline{Demostración.}

(a) En vista de la sucesión espectral
\[
 \check H^p\bigl(B|R,\underline H^q(F)\bigr)
 \Longrightarrow H^{p+q}(B|R,F)
\]
y de la igualdad
$\check H^0(\underline H^q(F))=0$ para $q>0$, basta demostrar que
$\check{\underline H}^{\,n}(G)$ es de tipo finito para todo $n$ y todo
prehaz $G$ de tipo finito. Por el Lema 3.2,
\[
 \check H^{\bullet}(B|R,G)=
 H^{\bullet}(\{P\longrightarrow B\},G),
\]
donde $P$ es un álgebra polinómica sobre $R$ y
$\{P\longrightarrow B\}\in\operatorname{Cov}(B)$. Este último es la
cohomología del complejo

%%%%%%%%%%  scan idx 86  |  folio 75  %%%%%%%%%%

\[
 \cdots
 G(P\times_BP\times_BP)
 \mathrel{\substack{\longleftarrow\\[-.6ex]\longleftarrow\\[-.6ex]\longleftarrow}}
 G(P\times_BP)
 \mathrel{\substack{\longleftarrow\\[-.6ex]\longleftarrow}}
 G(P).
\]
Si $B$ es una $R$-álgebra de tipo finito, puede elegirse $P$ de tipo
finito sobre $R$. Todos los productos fibrados iterados
$P\times_B\cdots\times_BP$ son entonces $R$-álgebras de tipo finito. Si
$G$ es de tipo finito, sus imágenes por $G$ son $A$-módulos de tipo
finito. Por consiguiente,
$H^{\bullet}(\{P\longrightarrow B\},G)$ es un $A$-módulo de tipo finito
porque $A$ es noetheriano.

(b) Sea $P$ un álgebra polinómica sobre $R$. Toda aplicación suprayectiva
de $R$-álgebras $B\longrightarrow P$ admite una sección, por lo que
\[
 H^i(\{B\longrightarrow P\},G)=0\qquad(i>0).
\]
Se sigue que $\check H^i(P|R,G)=0$ para todo $i>0$ y todo prehaz $G$
sobre $\underline C_{P|R}$. Por consiguiente, la sucesión espectral
\[
 \check H^p\bigl(P|R,\underline H^q_{P|R}(F)\bigr)
 \Longrightarrow H^{p+q}(P|R,F)
\]
degenera, y
\[
 \check H^0\bigl(P|R,\underline H^n(F)\bigr)
 \simeq H^n(P|R,F)
\]
para todo $n$. Sin embargo, para un haz $F$,
$\check H^0(P|R,\underline H^n(F))=0$ siempre que $n>0$. Así,
$H^n(P|R,F)=0$ para todo $n>0$.

(c) Si
$0\longrightarrow F'\longrightarrow F\longrightarrow F''
\longrightarrow0$ es exacta en
$\widetilde{\underline C}_{A|R}$, entonces la sucesión
\[
 0\longrightarrow F'(B)\longrightarrow F(B)
 \longrightarrow F''(B)
 \longrightarrow H^1(B|R,F')\longrightarrow H^1(B|R,F)
 \longrightarrow\cdots
\]
es exacta. Si $P\in\operatorname{Ob}(\underline C_{A|R})$ es un
álgebra polinómica sobre $R$, entonces $H^1(P|R,F')=0$ por (b), y por
tanto
\[
 0\longrightarrow F'(P)\longrightarrow F(P)
 \longrightarrow F''(P)\longrightarrow0
\]
es exacta.

Recíprocamente, sea $\underline P_{A|R}$ la subcategoría plena de
$\underline C_{A|R}$ con
\[
 \operatorname{Ob}(\underline P_{A|R})=
 \{P\in\operatorname{Ob}(\underline C_{A|R})\mid
   P\text{ es un álgebra polinómica sobre }R\},
\]
dotada de la topología inducida, y sea
$p:\underline P_{A|R}\longrightarrow\underline C_{A|R}$ la inclusión.
Todo objeto de $\underline C_{A|R}$ está recubierto por un objeto de
$\underline P_{A|R}$; el lema de comparación muestra, por tanto, que
\[
 \widetilde p:\widetilde{\underline C}_{A|R}
 \longrightarrow\widetilde{\underline P}_{A|R}
\]
es una equivalencia de categorías. Si la sucesión que se obtiene al
evaluar $F'\longrightarrow F\longrightarrow F''$ en cada
$P\in\operatorname{Ob}(\underline P_{A|R})$ es exacta, entonces
\[
 0\longrightarrow\widetilde pF'\longrightarrow\widetilde pF
 \longrightarrow\widetilde pF''\longrightarrow0
\]
es exacta en $\widetilde{\underline P}_{A|R}$ y, por consiguiente, la
sucesión original es exacta en
$\widetilde{\underline C}_{A|R}$. \hfill$\square$

%%%%%%%%%%  scan idx 87  |  folio 76  %%%%%%%%%%

\underline{Observación 3.4.} Toda flecha suprayectiva en
$\underline P_{A|R}$ admite una sección. Se sigue inmediatamente que
todo prehaz sobre $\underline P_{A|R}$ es un haz; es decir,
\[
 \widetilde{\underline P}_{A|R}=
 \overset{\vee}{\underline P}_{A|R}.
\]
Por consiguiente, el funtor de restricción
\[
 \widetilde{\underline C}_{A|R}
 \longrightarrow\overset{\vee}{\underline P}_{A|R}
\]
es una equivalencia de categorías.
